%% kapa: exam=1 round=22 year=2011 period=2 form=A subject=appraisal_law subject_ko=부동산관계법규 seq=4 number=44
%% source: Q-Net 감정평가사 기출문제 (cst00309) artlSeq=5036790 file=extracted/1008/5036790/2011년 제22회 감정평가사 1차시험 기출문제/2교시.hwp
%% conversion: kordoc 4.24.51 -> tools/gen1.py (auto); review=proofread-gpt-astra; flags=none
\begin{question}{44}
국토의 계획 및 이용에 관한 법령상 공공시설의 귀속에 관한 설명으로 옳지 않은 것은?
\begin{choices}
  \item 개발행위허가를 받은 자가 행정청인 경우 개발행위허가를 받은 자는 개발행위가 끝나 준공검사를 마친 때에는 해당 시설의 관리청에 공공시설의 종류와 토지의 세목을 통지하여야 한다.
  \item 개발행위허가를 받은 행정청이 새로 공공시설을 설치한 경우 당해 공공시설은 그 시설을 관리할 관리청에 무상으로 귀속된다.
  \item 개발행위허가를 받은 행정청에게 종래의 공공시설의 소유권이 무상으로 귀속된 경우 당해 행정청은 당해 공공시설의 처분으로 인한 수익금을 도시계획사업의 목적으로만 사용하여야 한다.
  \item 개발행위허가를 받은 자가 행정청이 아닌 경우 개발행위허가를 받은 자가 새로 설치한 공공시설은 그 시설을 관리할 관리청에 무상으로 귀속된다.
  \item 개발행위허가를 받은 자가 행정청이 아닌 경우 개발행위로 용도가 폐지되는 공공시설은 「국유재산법」과 「공유재산 및 물품 관리법」에도 불구하고 새로 설치한 공공시설의 설치비용에 상당하는 범위에서 개발행위허가를 받은 자에게 무상으로 양도되어야 한다.
\end{choices}
\end{question}
\answer{5}
